%-------------------------------------------------------------------------------
% book.tex  --  Designing Musical Networks
% Master file.  Chapters live in chapters/ (hand-written) and generated/
% (produced by pandoc from ../notes/*.md).  See book/Makefile.
%-------------------------------------------------------------------------------

\documentclass[
    fontsize=10pt,
    twoside=true,        % proper book layout: even/odd pages
    secnumdepth=2,       % number chapters, sections, subsections
    fontmethod=tex,      % Palatino via newpx + pdflatex.  Uses fonts
                         % that ship with every standard TeX Live install
                         % (no Libertinus, no Liberation Mono, no fontspec).
                         % The preprocess_note.py script normalises the
                         % notes' Unicode chars before pandoc, so pdflatex
                         % is sufficient.
]{kaobook}

% --- Language / quotes ---
\usepackage[english]{babel}
\usepackage[english=british]{csquotes}

% --- Math / theorems / refs ---
\usepackage{kaotheorems}
\usepackage{kaorefs}
\usepackage{epigraph}    % chapter-opening quotations

% --- Code listings -------------------------------------------------------------
% We use the listings package (portable, no shell-escape needed) with a
% C++-friendly style.  Code in the notes (markdown ```cpp ... ```) is
% converted by pandoc to lstlisting environments that pick this style up.
\usepackage{listings}
\usepackage{xcolor}

\definecolor{codebg}{HTML}{F7F7F4}
\definecolor{codekw}{HTML}{1F4E79}
\definecolor{codecm}{HTML}{6E7781}
\definecolor{codest}{HTML}{8B5E3C}

\lstdefinestyle{weft}{
    backgroundcolor   = \color{codebg},
    basicstyle        = \ttfamily\small,
    keywordstyle      = \color{codekw}\bfseries,
    commentstyle      = \color{codecm}\itshape,
    stringstyle       = \color{codest},
    numbers           = none,
    showstringspaces  = false,
    breaklines        = true,
    frame             = none,
    aboveskip         = 1.2em,
    belowskip         = 1.2em,
    xleftmargin       = 0.5em,
    xrightmargin      = 0.5em,
}
\lstset{style=weft}

% --- Bibliography --------------------------------------------------------------
\usepackage{kaobiblio}
\addbibresource{bibliography/refs.bib}

% --- Graphics paths ------------------------------------------------------------
\graphicspath{{images/}{./}}

% --- Custom commands shared across parts --------------------------------------
% These keep notation consistent between the geometric (Part I) and
% computational (Part II) chapters.
\newcommand{\R}{\mathbb{R}}
\newcommand{\C}{\mathbb{C}}
\newcommand{\N}{\mathbb{N}}
\newcommand{\E}{\mathbb{E}}
\newcommand{\inner}[2]{\langle #1, #2 \rangle}
\newcommand{\norm}[1]{\lVert #1 \rVert}
\newcommand{\abs}[1]{\lvert #1 \rvert}
\newcommand{\xb}{\mathbf{x}}
\newcommand{\yb}{\mathbf{y}}
\newcommand{\zb}{\mathbf{z}}
\newcommand{\wb}{\mathbf{w}}
\newcommand{\Wb}{\mathbf{W}}
% Convolution star reads better with a thin space on either side.
\newcommand{\conv}{\,*\,}

% --- Pandoc helper commands ----------------------------------------------------
% Pandoc's markdown -> latex conversion emits a few commands that aren't
% defined by LaTeX itself; pandoc's own default template provides them.
% Since we use a custom (minimal) template, we provide them here.
\providecommand{\tightlist}{%
    \setlength{\itemsep}{0pt}\setlength{\parskip}{0pt}}
% \passthrough wraps inline code (markdown backticks) so it survives
% pandoc's tokenization unchanged.  Default behaviour is just the content.
\providecommand{\passthrough}[1]{#1}

\makeindex[columns=3, title=Alphabetical Index, intoc]

%-------------------------------------------------------------------------------
\begin{document}

\titlehead{}
\subject{}
\title[Designing Musical Networks]{Designing Musical Networks}
\subtitle{From Linear Algebra to Music Composition}
\author[C.-E. Cella]{Carmine-Emanuele Cella}
\date{\today}
\publishers{draft --- not for circulation}

\frontmatter

\makeatletter
\uppertitleback{\@titlehead}
\lowertitleback{%
    \textbf{Draft.}\\
    This is a working draft of \emph{Designing Musical Networks}.  Not for
    circulation.  Comments welcome.

    \medskip

    \textbf{Colophon.}\\
    Typeset with \href{https://sourceforge.net/projects/koma-script/}{\KOMAScript{}} and \LaTeX{}
    using the \href{https://github.com/fmarotta/kaobook/}{kaobook} class.
    Code companion: \texttt{weft}, a header-only C++17 neural network library
    available alongside this book.
}
\makeatother

\dedication{%
    For those who hear mathematics in music\\
    and music in mathematics.\\
    \flushright -- the author
}

\maketitle

% --- Preface ---
\chapter*{Preface}
\input{chapters/preface.tex}

% --- TOC ---
\begingroup
\setlength{\textheight}{230\vscale}
\etocstandarddisplaystyle
\etocstandardlines
\tableofcontents
\endgroup

%-------------------------------------------------------------------------------
\mainmatter
\setchapterstyle{kao}

% =============================================================================
% PART I --- FOUNDATION
% =============================================================================
\pagelayout{wide}
\addpart{Foundation}
\pagelayout{margin}

% chapters/01_signals_as_vectors.tex
%
% Part I, Chapter 1.  Adapted and rewritten from "A Geometric Tour of
% Machine Learning" Chapter 1 (Mathematical Foundation), with a new
% musically-motivated opening, tighter prose, and margin notes to keep the
% reader oriented.

\chapter{Signals as vectors}
\labch{signals_as_vectors}

\epigraph{The fool doth think he is wise, but the wise man knows himself
to be a fool.}{William Shakespeare, \emph{As You Like It}}

%------------------------------------------------------------------------
\section{Two notes}
%------------------------------------------------------------------------

\noindent Listen to a sustained violin note and a sustained flute note,
played at the same pitch.  They have the same fundamental frequency.
They sound nothing alike.  Now play them together.  They blend into
something that is neither violin nor flute --- a third instrument that
exists only while they are sounding together.

\marginnote{%
This is the setup.  The math we develop in this chapter --- vector spaces,
inner products, projections --- is the language in which the difference
between a violin and a flute, and the meaning of ``blend,'' can be made
precise.}

The question \emph{why do they sound different} can be answered at two
levels.  At the musical level the answer is timbre: the way each
instrument's overtones reinforce or attenuate the others, the attack
transient, the bow noise or the breath noise, the small fluctuations in
pitch and amplitude.  At the mathematical level the answer is that the
two recordings, considered as long sequences of numbers, point in
different directions in a very high-dimensional space.  When we add them,
we are doing geometry: the result is a third point in the same space,
neither of the originals but related to both.

This book lives in the gap between those two answers.  The technical
chapters (Part~\textsc{II}) build software that operates on sounds as
high-dimensional vectors.  The creative chapters (Part~\textsc{III})
reason about timbre, style, and composition in those terms.  This first
part is about establishing the geometry --- the structure of the space
that makes the technical work possible and the creative work coherent.

The central claim is simple, and it will reappear in many guises:
\emph{representing a signal means projecting it onto a set of reference
vectors}.  The Fourier transform projects onto sinusoids.  The wavelet
transform projects onto localized oscillations.  Principal component
analysis projects onto data-derived directions of maximum variance.  A
neural network projects onto weight vectors learned from data, then
applies a nonlinearity and projects again.  The mathematics underlying
all four is the same: an inner product, repeated.  What differs is the
choice of reference vectors and whether they are designed by hand or
learned from data.

\marginnote{%
\textbf{Notation.}  Vectors are denoted by bold lowercase letters
$\xb$, $\yb$.  Matrices by bold uppercase, $\Wb$, $\mathbf{B}$.  The
$n$-th component of vector $\xb$ is written $\xb[n]$.  The inner product
of two vectors is $\inner{\xb}{\yb}$, and convolution is $\xb \conv \yb$.
Appendix~\ref{ch:matrix_theory} reviews matrix conventions.}

%------------------------------------------------------------------------
\section{Vector spaces}
%------------------------------------------------------------------------

To think of a signal as a point in space, we need to say precisely what
``space'' means.  A \emph{vector space} is a set $V$ on which two
operations are defined --- addition of elements, and multiplication by
scalars --- which satisfy a list of axioms.\sidenote{The full axiom list
appears in Appendix~\ref{ch:matrix_theory}; for our purposes, what matters
is that they let us add signals and scale them without leaving the space.}
The elements of $V$ are called \emph{vectors}.  When the scalar field is
the real numbers $\R$, we say $V$ is a real vector space; when it is the
complex numbers $\C$, a complex vector space.

A discrete-time signal of length $N$ is naturally an element of $\R^N$
(or $\C^N$, if we allow complex samples).  The $n$-th component is the
$n$-th sample.  Adding two signals samplewise is vector addition.
Multiplying every sample by a constant is scalar multiplication.  The
zero vector is the silent signal.

In this book we work exclusively in finite dimensions.  The infinite-dim\-en\-sion\-al
theory of Hilbert spaces is beautiful and important, but everything we
need can be done in $\R^N$ or $\C^N$, and the finite case is what the
computer manipulates anyway.

%------------------------------------------------------------------------
\section{Norm: the size of a signal}
%------------------------------------------------------------------------

To say one signal is louder than another, or that one is close to
another, we need a notion of size and of distance.  A \emph{norm} on a
vector space $V$ is a function $\norm{\cdot}: V \to \R$ satisfying
\emph{(i)} $\norm{\xb} \ge 0$ with equality only at $\xb = \mathbf{0}$,
\emph{(ii)} $\norm{\alpha \xb} = \abs{\alpha} \norm{\xb}$ for any scalar
$\alpha$, and \emph{(iii)} the triangle inequality $\norm{\xb + \yb} \le
\norm{\xb} + \norm{\yb}$.

The norm we will use almost everywhere is the \emph{Euclidean} or
\emph{$L^2$} norm,
\begin{equation}
    \norm{\xb}_2 \;=\; \sqrt{\sum_{n=0}^{N-1} \abs{\xb[n]}^2}.
    \label{eq:l2norm}
\end{equation}
The square of this quantity, $\norm{\xb}_2^2$, has a physical
interpretation when $\xb$ is an audio signal: it is the \emph{energy} of
the signal.  We will often work with energy directly rather than with
the norm, because it avoids the square root.

\marginnote{%
Other norms exist and matter.  The $L^1$ norm $\sum_n \abs{\xb[n]}$ is
used in sparse coding because it prefers solutions with many zeros.  The
$L^\infty$ norm $\max_n \abs{\xb[n]}$ is the peak amplitude.  We will
meet $L^1$ again in Part~\textsc{III} when we discuss matching pursuit.}

A vector space equipped with a norm is called a \emph{Banach space}.
Banach spaces are what we need to talk about size and convergence.  We
need one more piece of structure --- an inner product --- to talk about
angles and projections.  That promotes a Banach space to a Hilbert space.

%------------------------------------------------------------------------
\section{Inner product: global similarity}
%------------------------------------------------------------------------

The inner product is the single most important operation in this book.
For two real-valued signals $\xb$ and $\yb$ of length $N$,
\begin{equation}
    \inner{\xb}{\yb} \;=\; \sum_{n=0}^{N-1} \xb[n]\, \yb[n].
    \label{eq:inner}
\end{equation}
For complex-valued signals, the second factor is conjugated:
$\inner{\xb}{\yb} = \sum_n \xb[n]\, \overline{\yb[n]}$.

Geometrically, the inner product measures \emph{similarity}.  The clean
statement is the following.  Two vectors with the same magnitude have a
large inner product when they point in the same direction, zero when
they are perpendicular, and a large negative value when they point in
opposite directions.\sidenote{The relation
$\inner{\xb}{\yb} = \norm{\xb}\, \norm{\yb} \cos\theta$, where $\theta$
is the angle between them, makes this precise.  When $\theta = 0$ the
cosine is $1$; when $\theta = \pi/2$ it is zero; when $\theta = \pi$ it
is $-1$.}  When two vectors are identical, their inner product is the
square of their norm: $\inner{\xb}{\xb} = \norm{\xb}^2$.  When they are
orthogonal ($\inner{\xb}{\yb} = 0$), they share nothing along each
other's direction.

The inner product is \emph{global}.  Every sample of $\xb$ is multiplied
with every corresponding sample of $\yb$, and the products are summed
into a single number.  No spatial or temporal locality is preserved ---
the contribution of a sample at the beginning of the signal is
indistinguishable, in the final sum, from the contribution of a sample
at the end.  This is sometimes what we want (when comparing the overall
frequency content of two signals, say) and sometimes not (when looking
for a specific pattern at an unknown time).  When we want locality, we
use convolution instead, which is the next section's subject.

A vector space equipped with an inner product (and the norm derived
from it via $\norm{\xb} = \sqrt{\inner{\xb}{\xb}}$) is called a
\emph{Hilbert space}.  Hilbert spaces are the natural setting for
analysis, synthesis, and learning.  Every space we work with in this book
will be a Hilbert space.

%------------------------------------------------------------------------
\section{Convolution: local similarity}
\label{sec:convolution}
%------------------------------------------------------------------------

A signal is rarely interesting in its entirety.  When we listen, we
attend to local features --- the attack of a note, a particular harmonic
fluctuation, a pause.  The mathematical operation that picks out local
features is the \emph{convolution}.

Given a signal $\xb$ and a short pattern $\mathbf{h}$ (called the
\emph{kernel} or \emph{filter}) of length $M$, the convolution
$\xb \conv \mathbf{h}$ is the sequence
\begin{equation}
    (\xb \conv \mathbf{h})[n]
    \;=\;
    \sum_{m=0}^{M-1} \xb[n+m]\, \mathbf{h}[m].
    \label{eq:conv}
\end{equation}
For each output position $n$, the kernel is aligned with the signal
starting at sample $n$, and the inner product is taken over just those
$M$ samples.  This is the local version of the inner product from the
previous section: instead of computing one number that summarises the
entire signal, we compute many numbers, each summarising a small
neighbourhood.

\marginnote{%
A small example.  If $\mathbf{h} = (1, -1)$, then
$(\xb \conv \mathbf{h})[n] = \xb[n+1] - \xb[n]$, the discrete derivative.
If $\mathbf{h} = (1, 1)/2$, we get a two-point smoother.  Choosing
$\mathbf{h}$ chooses what features the convolution highlights.  This is
precisely what a convolutional neural network learns to do: its kernels
are exactly such patterns, but discovered from data rather than designed.}

Convolution sits at the centre of both signal processing and modern
machine learning.  In signal processing it underlies filtering,
reverberation, pattern detection, and edge detection.  In machine
learning it is the basis of convolutional neural networks, which we
build in Part~\textsc{II}.  In both fields, convolution is doing the
same thing: \emph{repeated inner products, slid along the signal}.

The kernel $\mathbf{h}$ is itself a signal --- a short one, typically.
Choosing the kernel is choosing what local feature to look for.  A
kernel that resembles a quarter-period of a sinusoid will respond
strongly where the signal looks like that fragment of sinusoid; this is
the seed of Fourier analysis.  A kernel that resembles an edge will
respond at sharp transitions; this is the seed of the Sobel filter and
of the first layer of an image conv-net.  Different kernels reveal
different features; the right kernel to use depends on what we are
trying to discover.

There is a deep theorem connecting convolution to Fourier analysis,
namely that convolution in the time domain is equivalent to pointwise
multiplication in the frequency domain.  We will state and prove it in
Chapter~\ref{ch:apriori} when we have set up the Fourier transform;
Appendix~\ref{ch:convolution_theorem} contains a fuller proof.  The
practical implication is that convolution can be computed quickly via
the FFT, which is why convolutional networks scale.

%------------------------------------------------------------------------
\section{Basis and projection}
%------------------------------------------------------------------------

We have inner products (global similarity) and convolutions (local
similarity).  The next step is to use them to \emph{decompose} a signal
into pieces.

Given a set of vectors $\{\mathbf{b}_1, \mathbf{b}_2, \ldots,
\mathbf{b}_K\}$ in a vector space $V$, the \emph{span} of the set is the
collection of all linear combinations
\(
\sum_k \alpha_k \mathbf{b}_k
\)
that can be formed from them.  If every vector in $V$ can be written
this way --- if the span of the set is all of $V$ --- and if no vector
in the set is a linear combination of the others, the set is called a
\emph{basis} for $V$.

A basis is the minimal set of building blocks for the space.  In $\R^N$
there are always exactly $N$ vectors in any basis; we can pick the
``natural'' basis $\mathbf{e}_n$ (the $n$-th unit vector, equal to one
at position $n$ and zero elsewhere) or any other set of $N$ linearly
independent vectors.

The point of having a basis is that any vector $\xb \in V$ can be
written uniquely as
\begin{equation}
    \xb \;=\; \sum_{k=1}^{N} c_k\, \mathbf{b}_k.
    \label{eq:expansion}
\end{equation}
The numbers $c_k$ are the \emph{coordinates} of $\xb$ in the basis
$\{\mathbf{b}_k\}$.  They are what change when we change basis; the
underlying vector $\xb$ does not.

When the basis is \emph{orthonormal} --- meaning the basis vectors are
mutually perpendicular and each has unit norm --- there is a beautifully
simple formula for the coordinates:
\begin{equation}
    c_k \;=\; \inner{\xb}{\mathbf{b}_k}.
    \label{eq:coord_orth}
\end{equation}
The coordinate $c_k$ is just the inner product of $\xb$ with the $k$-th
basis vector.  Combining this with~\eqref{eq:expansion},
\begin{equation}
    \xb \;=\; \sum_{k=1}^{N} \inner{\xb}{\mathbf{b}_k}\, \mathbf{b}_k.
    \label{eq:reconstruction}
\end{equation}
This is the single most important formula in classical signal processing.
It says: to reconstruct $\xb$, compute its inner product with each basis
vector (the \emph{analysis} step), then build $\xb$ back up by combining
the basis vectors weighted by those inner products (the \emph{synthesis}
step).  In an orthonormal basis, analysis and synthesis are exact
inverses of one another.

\marginnote{%
When the basis is not orthonormal, equation~\eqref{eq:coord_orth} still
makes sense as the \emph{projection} of $\xb$ onto the direction of
$\mathbf{b}_k$, but the projections no longer give exact coordinates ---
the basis vectors interfere.  An orthonormal basis decouples them.
This is why so many useful representations (Fourier, the standard
basis, principal components) are designed to be orthonormal.}

We will see many examples of \eqref{eq:reconstruction} in disguise.
The Fourier transform is this formula with $\mathbf{b}_k$ a sinusoid of
frequency $k$.  The Discrete Wavelet Transform is this formula with
$\mathbf{b}_k$ a scaled, shifted copy of a mother wavelet.  Principal
component analysis is this formula with $\mathbf{b}_k$ the $k$-th
eigenvector of the data covariance.  A neural network with linear
activations is this formula with $\mathbf{b}_k$ a learned weight vector.

\emph{The choice of basis is the design of the representation.}  When we
pick sinusoids, we have decided that frequency is the property we want
to expose.  When we pick wavelets, we have decided that local
time-frequency structure matters.  When we let a network learn its
basis, we have decided to let the data decide.  Each choice has
consequences: it determines what the representation makes obvious, what
it hides, and what kinds of operations are easy versus hard to express
in it.  Part~\textsc{II} of this book builds the tools to make such
choices; Part~\textsc{III} explores their compositional consequences.

%------------------------------------------------------------------------
\section{Analysis and synthesis as matrices}
%------------------------------------------------------------------------

Equations~\eqref{eq:coord_orth} and~\eqref{eq:reconstruction} both
become cleaner when written in matrix form.  Let $\mathbf{B}$ be the
matrix whose $k$-th \emph{row} is the basis vector $\mathbf{b}_k^\top$.
Then computing all the inner products $\inner{\xb}{\mathbf{b}_k}$ at
once is the same as multiplying $\xb$ by $\mathbf{B}$:
\begin{equation}
    \mathbf{c} \;=\; \mathbf{B}\, \xb \qquad \text{(analysis)}.
    \label{eq:analysis_matrix}
\end{equation}
The vector $\mathbf{c}$ contains all the coordinates $c_k$.  Equivalently
written, $\mathbf{c} = \Phi(\xb)$ where $\Phi$ is the \emph{analysis
operator}.

For the synthesis direction, if the basis is orthonormal, the inverse of
$\mathbf{B}$ is its transpose (or, in the complex case, its Hermitian
adjoint $\mathbf{B}^H$).  Multiplying~\eqref{eq:analysis_matrix} on the
left by $\mathbf{B}^H$ recovers $\xb$:
\begin{equation}
    \xb \;=\; \mathbf{B}^H \mathbf{c} \;=\; \mathbf{B}^H\, \mathbf{B}\, \xb
    \qquad \text{(synthesis)}.
    \label{eq:synthesis_matrix}
\end{equation}
For non-orthonormal but still invertible $\mathbf{B}$, the synthesis
uses $\mathbf{B}^{-1}$ rather than $\mathbf{B}^H$.

So the entire analysis-synthesis cycle is just a pair of matrix-vector
multiplications.  This is the bridge to the rest of the book.  The
Fourier transform is exactly this with one specific choice of
$\mathbf{B}$ (a matrix of complex sinusoids).  The DWT is this with a
different specific choice (a matrix of wavelets at multiple scales).  A
linear neural-network layer is this with $\mathbf{B}$ \emph{learned from
data} instead of designed in advance.

%------------------------------------------------------------------------
\section{From basis to representation}
\label{sec:representations}
%------------------------------------------------------------------------

The basis idea generalises in three important ways, and the
generalisations are what make modern representations work.

\paragraph{Relaxing linear independence.}
In a classical basis, we insist that no basis vector is a linear
combination of the others.  This is what makes the representation
\emph{unique}: every $\xb$ has exactly one set of coordinates.  But
uniqueness is not always what we want.  Sometimes we want a redundant
set of vectors --- more than $N$ of them in an $N$-dimensional space ---
because the redundancy can encode information about the signal that a
minimal basis cannot.  Such sets are called \emph{frames} or
\emph{dictionaries}.  Matching pursuit, which we will meet in
Part~\textsc{III}, is a method for choosing a sparse subset of a
dictionary to represent a signal.

\paragraph{Allowing dimensionality change.}
In a classical basis, the matrix $\mathbf{B}$ is square: $N$ basis
vectors in an $N$-dimensional space.  If we allow $\mathbf{B}$ to be
rectangular --- say $K \times N$ with $K < N$ --- we get a
\emph{projection onto a subspace}.  The output has fewer dimensions than
the input.  This is dimensionality reduction, and it is what PCA, the
encoder of an autoencoder, and the early layers of a classifier all do.
The inverse direction --- $N$ to $K$ with $K > N$ --- is dimensionality
expansion, which is what the decoder of an autoencoder does, and what
generative models do when they turn a small latent vector into a full
image or audio file.

\paragraph{Nonlinearity.}
A pure projection $\xb \mapsto \mathbf{B}\xb$ is a linear operation;
composing two linear operations gives another linear operation, so
\emph{stacking} pure projections gives nothing new --- any depth of
linear layers collapses to a single equivalent layer.  To get a deep
representation that is more expressive than a shallow one, we need to
insert a nonlinearity between the projections.  In modern neural
networks, this is the ReLU function $\max(0, \cdot)$, the sigmoid
$\sigma$, or any of a small zoo of others.  In the Scattering Transform
(Chapter~\ref{ch:apriori}), it is the absolute value $\abs{\cdot}$ that
plays this role.

Putting all three together --- redundant dictionaries, varying
dimensionality, nonlinear stages --- gives the basic shape of a modern
representation:
\begin{equation}
    \mathbf{z}^{(k)} \;=\; f_k\!\left(\mathbf{B}_k\, \mathbf{z}^{(k-1)} + \mathbf{b}_k\right),
    \qquad k = 1, 2, \ldots, L,
    \label{eq:layered}
\end{equation}
with $\mathbf{z}^{(0)} = \xb$ the input, $\mathbf{B}_k$ the projection
matrix of the $k$-th layer, $\mathbf{b}_k$ a bias vector (allowing the
layer to be \emph{affine} rather than purely linear, which lets the
representation work with non-zero-centred data), and $f_k$ the
nonlinearity.

\marginnote{%
\textbf{Encoders and decoders.}  When the layers reduce dimensionality
($\mathbf{B}_k$ progressively narrower) and end in a small latent
vector, we call the whole stack an \emph{encoder}.  Reverse it ---
layers that progressively expand dimensionality from a latent vector
back to a full signal --- and we have a \emph{decoder}.  An
autoencoder is an encoder followed by a decoder, trained so that the
output reconstructs the input.}

Equation~\eqref{eq:layered} is the structural backbone of essentially
every neural network in this book.  Chapter~\ref{ch:learned} discusses
how the $\mathbf{B}_k$ and $\mathbf{b}_k$ are \emph{learned} from data
by gradient descent on a loss function.  In Part~\textsc{II} we will
build the code that does this.

%------------------------------------------------------------------------
\section{What we have so far, and what comes next}
%------------------------------------------------------------------------

We have set up the conceptual frame:

\begin{itemize}
    \item Signals are vectors in a finite-dimensional Hilbert space.
    \item The inner product measures global similarity between vectors.
    \item Convolution is the inner product applied locally, kernel-by-kernel.
    \item Representations are projections onto a chosen basis (or
    dictionary), and reconstruction is the reverse operation.
    \item The choice of basis is the design of the representation; it
    determines what is made obvious and what is hidden.
    \item Layered representations stack projections separated by
    nonlinearities, allowing them to express far more than a single linear
    projection can.
\end{itemize}

\noindent The next two chapters give specific instances.
Chapter~\ref{ch:apriori} treats \emph{a-priori} representations --- those
whose basis is designed in advance --- with the Fourier transform, the
wavelet transform, and the Scattering Transform as our three main
examples.  Chapter~\ref{ch:learned} treats \emph{learned} representations
--- those whose basis is derived from data --- with principal component
analysis, $k$-means clustering, matching pursuit, and neural networks.
Chapter~\ref{ch:style_transfer} closes Part~\textsc{I} by working a
single concrete problem --- musical style transfer --- through each of
these representations in turn, showing how the choice changes what is
possible.  This sets up the technical work of Part~\textsc{II}, where we
build the code, and the compositional work of Part~\textsc{III}, where
we use it.

% chapters/02_apriori_representations.tex
%
% Part I, Chapter 2.  A-priori representations: Fourier, Wavelet,
% Scattering.  Adapted from the geometric tour's Chapter 2 with the
% Scattering Transform significantly expanded (it's the bridge to
% Chapter 3) and the connection to convolution made more explicit.

\chapter{A-priori representations}
\labch{apriori}

\epigraph{Music is the pleasure the human mind experiences from counting
without being aware that it is counting.}{Gottfried Wilhelm Leibniz}

%------------------------------------------------------------------------
\section{The question and three answers}
%------------------------------------------------------------------------

\noindent A sustained violin note arriving at our ears is a pressure
wave: a long sequence of small fluctuations around atmospheric pressure,
sampled by the eardrum at something like 40\,000 times per second.  As a
mathematical object it is a vector in $\R^N$ with $N$ very large.  As a
listening experience, it is a single sustained tone with a particular
timbre, attack, and decay.

The gap between those two descriptions is what a \emph{representation} is
supposed to close.  We want a transformation $\Phi$ such that $\Phi(\xb)$
exposes the structure we care about --- pitch, brightness, attack, the
texture of the bowing --- and ideally hides what we don't.  The previous
chapter argued that any reasonable $\Phi$ looks like an inner product
with a chosen set of reference vectors.  This chapter looks at three
specific choices, each historically and practically important.  All
three are \emph{a-priori} in the sense that the reference vectors are
designed mathematically, not learned from data.  The next chapter will
turn to representations whose reference vectors come from data.

The three a-priori representations we consider are:

\begin{itemize}
    \item The \emph{Discrete Fourier Transform} (Section~\ref{sec:dft}),
    whose reference vectors are pure sinusoids.  Fourier exposes
    frequency content with maximum sharpness and is the workhorse of
    classical signal processing.

    \item The \emph{Discrete Wavelet Transform} (Section~\ref{sec:dwt}),
    whose reference vectors are localized oscillations --- a single
    \emph{mother wavelet}, scaled and translated to cover the
    time-frequency plane.  Wavelets give us frequency information with
    time localisation, at the cost of some spectral precision.

    \item The \emph{Scattering Transform} (Section~\ref{sec:scattering}),
    a multi-layer cascade of wavelet convolutions and absolute-value
    nonlinearities.  It is the bridge between classical signal processing
    and modern convolutional neural networks: structurally identical to a
    CNN, but with all the filters fixed by mathematics rather than learned
    from data.
\end{itemize}

\noindent The progression Fourier $\to$ Wavelet $\to$ Scattering will
turn out to add, one step at a time, three things: \emph{locality}
(wavelets are localised, sinusoids are not), \emph{depth} (Scattering
stacks two or three wavelet transforms), and \emph{nonlinearity}
(Scattering inserts an absolute value between the layers).  Once we have
these three, a fourth step --- replacing the predefined filters with
\emph{learned} ones --- gets us a convolutional network.  That fourth
step is the subject of the next chapter.

%------------------------------------------------------------------------
\section{The Discrete Fourier Transform}
\label{sec:dft}
%------------------------------------------------------------------------

The Discrete Fourier Transform represents a length-$N$ signal as a sum
of $N$ complex sinusoids.  For each frequency index $k \in
\{0, 1, \ldots, N-1\}$,
\begin{equation}
    \hat{\xb}[k]
    \;=\;
    \sum_{n=0}^{N-1} \xb[n] \, e^{-i 2\pi k n / N}.
    \label{eq:dft}
\end{equation}
The inverse transform reconstructs $\xb$ from $\hat\xb$ via the same
formula with the opposite sign in the exponent and a $1/N$ normalisation:
\begin{equation}
    \xb[n]
    \;=\;
    \frac{1}{N} \sum_{k=0}^{N-1} \hat{\xb}[k] \, e^{i 2 \pi k n / N}.
    \label{eq:idft}
\end{equation}

\marginnote{%
The DFT in \eqref{eq:dft} would naively cost $O(N^2)$ operations to
evaluate.  The \emph{Fast Fourier Transform} reduces this to
$O(N \log N)$ by exploiting the recursive structure of the sum.  The FFT
is what makes Fourier analysis practical for audio; we use it
extensively in Part~II.}

\paragraph{The DFT as projection.}  Compare \eqref{eq:dft} with the
inner-product formula \eqref{eq:inner} from the previous chapter.  The
DFT is just the inner product of $\xb$ with the $k$-th complex sinusoid:
\begin{equation}
    \hat{\xb}[k] \;=\; \inner{\xb}{\mathbf{f}_k},
    \qquad
    \mathbf{f}_k[n] \;=\; e^{i 2 \pi k n / N}.
    \label{eq:dft_as_inner}
\end{equation}
Each $\mathbf{f}_k$ is one of our reference vectors: a discrete sinusoid
that completes exactly $k$ full cycles over the length of the signal.
The set $\{\mathbf{f}_k\}_{k=0}^{N-1}$ is an \emph{orthogonal basis} for
$\C^N$, and in fact an orthonormal one after the $1/N$ scaling in
\eqref{eq:idft}.  This is why the inverse transform exists in closed
form: an orthogonal basis is its own decoder.

In matrix form, letting $\mathbf{F}$ be the $N \times N$ matrix with
$\mathbf{F}_{kn} = e^{-i 2 \pi k n / N}$,
\begin{equation}
    \hat\xb \;=\; \mathbf{F}\, \xb,
    \qquad
    \xb \;=\; \frac{1}{N} \mathbf{F}^H\, \hat\xb.
    \label{eq:dft_matrix}
\end{equation}
The Fourier transform is, exactly, a particular choice of
$\mathbf{B}$ in the framework of the previous chapter.

\paragraph{What the DFT exposes.}  Because each basis vector
$\mathbf{f}_k$ is a pure tone of frequency $k$, the coefficient
$\hat\xb[k]$ is large precisely when $\xb$ contains energy at that
frequency.  The magnitude $\abs{\hat\xb[k]}$ is the amount of frequency
$k$ in the signal; the argument $\arg \hat\xb[k]$ is its phase.  For
musical signals --- a sustained pitched tone --- the DFT magnitude
shows a peak at the fundamental and at each harmonic, telling us
immediately how loud each harmonic is.  This is the basis of pitch
detection, spectral envelopes, and timbre analysis.

\paragraph{What the DFT hides.}  Two properties of \eqref{eq:dft_as_inner}
matter.  First, the sum runs over the \emph{entire} signal: every sample
contributes to every coefficient.  A short click at $n = 100$
contributes to $\hat\xb[k]$ for every $k$.  The DFT has no notion of
``when'' --- it gives a global summary of frequency content, averaged
over all of time.  Second, the sinusoids $\mathbf{f}_k$ extend
unattenuated across the whole signal length; they are global functions
masquerading as basis vectors.

This is fine for stationary signals (an organ note holding constant for
seconds at a time) but disastrous for transient signals (a snare hit,
a consonant in speech, the attack of a piano note).  For such signals
we want \emph{locality}: a representation that tells us not just what
frequencies are present but \emph{when}.

\paragraph{Translation invariance and frequency sensitivity.}
The DFT has one beautiful invariance and one painful sensitivity.

If we shift $\xb$ in time by $\Delta$, the DFT magnitudes are
unchanged --- only the phases change:
\begin{equation}
    \xb[n - \Delta] \;\;\longleftrightarrow\;\;
    \hat\xb[k] \, e^{-i 2 \pi k \Delta / N}.
\end{equation}
This is what we want: a violin note arriving slightly later sounds the
same.  The magnitude spectrum is \emph{translation-invariant}.

But if we shift $\xb$ in \emph{frequency} by even a small amount --- say,
the violin is slightly mistuned, or our sample rate drifts ---
the DFT coefficients can move around dramatically.  A pure tone at
frequency $k + 0.5$ smears across many bins because it doesn't match
any single basis vector.  This is the well-known \emph{leakage}
phenomenon, and it's the fundamental cost of using strictly periodic
basis vectors on signals that are not strictly periodic.

\marginnote{%
A continuous version of the leakage problem is the
\emph{Heisenberg uncertainty principle for signals}: a representation
cannot have arbitrarily sharp resolution in both time and frequency.
The DFT chooses maximum frequency resolution (sharp basis vectors in
frequency) and accepts zero time resolution (global basis vectors in
time).  Wavelets choose a different trade-off; the scattering transform
yet another.}

\paragraph{Convolution becomes multiplication.}  The Fourier transform's
most useful algebraic property, far beyond its role as a representation,
is that convolution in the time domain corresponds to pointwise
multiplication in the frequency domain:
\begin{equation}
    \yb \;=\; \xb \conv \mathbf{h}
    \qquad \Longleftrightarrow \qquad
    \hat\yb[k] \;=\; \hat\xb[k] \cdot \hat{\mathbf{h}}[k].
    \label{eq:conv_theorem}
\end{equation}
This is the \emph{convolution theorem}; we sketch a proof in
Appendix~\ref{ch:convolution_theorem}.  The practical consequence is
that a convolution that would naively cost $O(NM)$ to evaluate directly
costs $O(N \log N)$ via three FFTs (forward, multiply, inverse).  For
large kernels, the FFT-based convolution is by far the faster method.

In the other direction, multiplication of two signals corresponds to
\emph{convolution} of their spectra.  This is how the Short-Time
Fourier Transform works (next), why amplitude modulation in audio
creates sidebands, and why the spectral analysis of a windowed signal
shows the window's spectrum convolved with the signal's spectrum.

%------------------------------------------------------------------------
\section{The Discrete Wavelet Transform}
\label{sec:dwt}
%------------------------------------------------------------------------

The DFT gives us frequency at the cost of time.  The Wavelet Transform
restores time by giving up some frequency.

\paragraph{Mother wavelets and their family.}  Pick a \emph{mother
wavelet} $\psi$: a short, oscillatory, zero-mean function of finite
energy.  Generate a family of basis vectors by scaling and translating
it:
\begin{equation}
    \psi_{a, b}[n]
    \;=\;
    \frac{1}{\sqrt{|a|}} \, \psi\!\left( \frac{n - b}{a} \right).
    \label{eq:wavelet_family}
\end{equation}
The scale parameter $a$ stretches or compresses the wavelet: a large
$a$ produces a wide, low-frequency wavelet; a small $a$ produces a
narrow, high-frequency one.  The translation parameter $b$ shifts the
wavelet along the time axis.  The normalisation $1/\sqrt{|a|}$ keeps the
energy constant across scales.

The Discrete Wavelet Transform of a signal $\xb$ is the array of inner
products of $\xb$ with this family:
\begin{equation}
    \mathbf{w}[a, b] \;=\; \inner{\xb}{\psi_{a, b}}.
\end{equation}
The coefficient $\mathbf{w}[a, b]$ tells us how much of the signal looks
like a wavelet of scale $a$ centred at time $b$.  Large coefficients
mark places in time-frequency where the signal has structure at that
particular scale.

\marginnote{%
Common mother wavelets include the \emph{Haar} (a square wave: simplest,
discontinuous), \emph{Daubechies} families (compactly supported,
orthogonal, smoother), the \emph{Mexican hat} (a Gaussian's second
derivative), and the \emph{Morlet} (a complex sinusoid times a Gaussian,
much used in time-frequency audio analysis).  The choice of mother
wavelet matters in practice; for this book the Morlet is the most
relevant.}

\paragraph{The time-frequency plane.}  Where the DFT gives one
coefficient per frequency, the DWT gives one coefficient per
\emph{location in time-frequency}.  Plot a DWT as a heatmap with time
on the horizontal axis and scale on the vertical axis, and you see a
picture: which frequencies are active when.  This is exactly the
\emph{spectrogram-like} visualisation that audio engineers use to look
at sounds.  In Part~II we build a related construction --- the Short-Time
Fourier Transform --- which gives a similar picture using DFTs of
short windowed segments.  Wavelets are the more principled relative,
with proper mathematical foundations rather than ad-hoc windowing.

\paragraph{What the wavelet representation costs.}  We gain time
localisation, but we pay for it.  Two specific costs.

First, the representation is no longer translation-invariant in
magnitude.  Shifting $\xb$ in time \emph{does} change which wavelet
coefficients are large, because the coefficients are indexed by both
scale and position.  A wavelet representation is translation-\emph{equivariant}:
the pattern of coefficients shifts in lockstep with the signal, but the
shifted pattern is not the same pattern.

Second, the frequency resolution at any single scale is lower than the
DFT's.  Each wavelet covers a range of frequencies rather than a single
one; the trade is wider time coverage in exchange for narrower
frequency coverage.  This is unavoidable: tight in time forces broad in
frequency, and vice versa.

\paragraph{Stability to frequency drift.}  In exchange, the DWT is much
more stable than the DFT to small drifts in frequency.  A pure tone at
frequency $f_0 + \delta$ produces wavelet coefficients close to those
of a pure tone at $f_0$, because each wavelet's frequency response is
already broad.  This makes wavelets the better tool for analysing
real-world signals whose frequencies wobble or evolve --- vibrato,
glissando, the slowly drifting pitch of a struck bell.

\paragraph{Matrix form.}  Just as with the DFT, the DWT has a matrix
representation.  Let $\boldsymbol\Psi$ be the matrix whose rows are the
wavelets $\psi_{a, b}$ at every $(a, b)$ pair in our chosen
discretisation.  Then
\begin{equation}
    \mathbf{w} \;=\; \boldsymbol\Psi\, \xb,
    \qquad
    \xb \;=\; \boldsymbol\Psi^H\, \mathbf{w}
    \quad\text{(if the wavelets form an orthonormal basis).}
    \label{eq:dwt_matrix}
\end{equation}
We have the same analysis-synthesis structure as the DFT, with a
different choice of basis.  The framework is invariant; the basis is
the design choice.

%------------------------------------------------------------------------
\section{Between transforms: the time-frequency trade}
%------------------------------------------------------------------------

Before we get to the Scattering Transform, it is worth pausing on the
relationship between the DFT and the DWT, because it makes the design
space visible.

The DFT picks one extreme of the time-frequency trade: maximally sharp
in frequency, maximally blurry in time.  Every basis vector spans the
whole signal, so we know exactly which frequencies are present but
nothing about when.

The Short-Time Fourier Transform (STFT), which we build in Part~II,
picks a middle point: we cut the signal into short overlapping
windows, take the DFT of each, and stack the results.  Each window has
fixed time width (call it $W$), so every frequency bin has the same
time resolution $W$ and the same frequency resolution $\sim\!1/W$.
Audio engineers reach for the STFT constantly; it is the basis of the
spectrogram.

The DWT picks a different middle point, one where the time-frequency
resolution varies with scale.  High frequencies (small $a$) get fine
time resolution and coarse frequency resolution --- a clear ``the
transient happened here'' answer.  Low frequencies (large $a$) get
coarse time resolution and fine frequency resolution --- a clear
``the pitch is $f_0$'' answer over a longer window.  This is exactly
the right behaviour for natural signals: short events tend to be
broadband, long events tend to be narrowband.

\marginnote{%
The DWT's variable resolution is sometimes called \emph{constant-Q
analysis}: each band has a fixed ratio of frequency width to centre
frequency.  This matches the way the human auditory system processes
sound, where critical bands also have approximately constant $Q$.  It is
not a coincidence that wavelet-like decompositions appear in the
cochlea.}

The progression is from rigid (DFT) to adaptive (DWT) to multi-layered
(Scattering, next).  Each step gives up some elegance --- the DFT has a
closed-form inverse; the DWT requires careful choice of mother wavelet;
the Scattering Transform has multiple stages and a fixed but non-trivial
schedule of operations --- in exchange for better behaviour on real
signals.

%------------------------------------------------------------------------
\section{The Scattering Transform}
\label{sec:scattering}
%------------------------------------------------------------------------

The Scattering Transform, introduced by Mallat in the early 2010s
(\emph{Group invariant scattering}, Communications on Pure and Applied
Mathematics, 2012), builds on the DWT in two crucial ways: it adds a
nonlinearity (the absolute value) between successive wavelet
decompositions, and it stacks several of them.  The resulting structure
is, mathematically and conceptually, a convolutional neural network with
all the filters fixed by hand.

\paragraph{One layer.}  We start with a signal $\xb$ and a family of
wavelets $\{\psi_{\lambda_1}\}$ indexed by scale (we collapse the
notation $\psi_{a, b}$ into $\psi_\lambda$ where $\lambda$ ranges over
both scale and shift).  The first-layer scattering coefficients are
\begin{equation}
    S_1 \xb \;=\; \abs{\xb \conv \psi_{\lambda_1}}.
    \label{eq:scattering1}
\end{equation}
Compute the wavelet transform of $\xb$, take the absolute value of each
coefficient.  This is the only operation; nothing learned, nothing
adjustable.

\paragraph{Two layers.}  The output $S_1 \xb$ is itself a signal --- in
fact, a collection of signals, one per scale.  We can apply the same
operation to each of them, using a second family of wavelets
$\{\psi_{\lambda_2}\}$:
\begin{equation}
    S_2 \xb \;=\; \abs{\, \abs{\xb \conv \psi_{\lambda_1}} \conv \psi_{\lambda_2} \,}.
    \label{eq:scattering2}
\end{equation}
This is a wavelet transform applied to the magnitude of a wavelet
transform.  The result captures \emph{interactions} between frequency
bands: how does activity at scale $\lambda_1$ vary over time?  How does
that variation itself look at scale $\lambda_2$?

The transform can be extended to three or more layers by repeating the
recipe.  In practice, two or three layers suffice; later layers tend
to capture diminishing amounts of energy.

\paragraph{Why the absolute value.}  Without the nonlinearity between
the two wavelet transforms, equation \eqref{eq:scattering2} would
collapse into a single wavelet transform (because the composition of
two linear operators is a linear operator).  The absolute value
prevents this collapse: it is the \emph{nonlinearity} that makes a deep
representation more expressive than a shallow one.

The choice of absolute value, specifically, is not arbitrary.  It is
the simplest nonlinearity that is \emph{stable} (Lipschitz: small
changes in input produce only small changes in output) and that
removes the phase information (which is highly oscillatory in time and
would otherwise destabilise the second layer).  Other nonlinearities
work too, but the absolute value has the cleanest theoretical
properties.

\paragraph{Stability and invariance.}  The scattering transform inherits
two key properties from its construction:

\begin{itemize}
    \item It is \emph{stable to deformations}: if $\xb$ is replaced by
    a slightly time-warped version $\xb \circ \tau$ where $\tau$ is a
    smooth deformation, the scattering coefficients $S\xb$ change by an
    amount controlled by the smoothness of $\tau$, not by the size of
    the deformation itself.  This is the right notion of robustness for
    natural signals, which are constantly subject to small distortions
    (recording-room acoustics, microphone characteristics, breath
    fluctuations).

    \item It is \emph{translation-invariant at large scale}: when the
    final layer is averaged over a window much wider than the signal's
    features, the result depends on $\xb$'s content but not on where in
    time that content occurs.  This is what we want for classification
    (a violin is a violin whether it arrives at sample $1000$ or
    $2000$).
\end{itemize}

\marginnote{%
\textbf{Why this matters.}  Bruna and Mallat showed in 2013 (\emph{Invariant
scattering convolution networks}, IEEE TPAMI) that a two-layer scattering
transform achieves near-state-of-the-art accuracy on some image and audio
classification tasks, \emph{without any learning}.  The features it
produces are good enough that a simple linear classifier on top reaches
accuracies that a CNN of the same depth would achieve.  This is a
remarkable result: it suggests that much of what early CNN layers do can
be replicated by fixed mathematical operations.  Of course, deeper
learned networks eventually outperform scattering, but the comparison
frames the question: what is the depth of learning truly buying us?}

\paragraph{The architecture.}  Look at the structure of the scattering
transform from a distance.  Each layer convolves the input with a bank
of filters, applies a nonlinearity, and passes the result to the next
layer.  Strip the words ``wavelet'' and ``absolute value'' and replace
them with ``learned filter'' and ``ReLU'' and you have, exactly, a
convolutional neural network.  Wavelet convolution becomes
$\Wb \conv \xb$; the absolute value becomes ReLU; the cascade of
layers is the same.

This is the bridge from classical signal processing to modern deep
learning.  The scattering transform is the \emph{minimum} deep network:
filters chosen for theoretical optimality (translation invariance,
locality, stability), nonlinearity chosen for stability.  A CNN
\emph{replaces} the principled-but-rigid wavelets with filters that the
data adjusts, trading mathematical guarantees for flexibility.  Some
problems do not benefit from that trade --- and for those, scattering is
competitive.  Most problems do --- and for those, we need learning.

\paragraph{Hierarchical features.}  The scattering transform also makes
explicit a structural point about \emph{depth}.  The first-layer
coefficients $\abs{\xb \conv \psi_{\lambda_1}}$ are very similar to a
spectrogram: they tell us about frequency content over time.  The
second-layer coefficients capture how that frequency content is itself
modulated.  Energy that fluctuates rapidly at the first layer becomes
visible as a clear peak at a particular $\lambda_2$ in the second layer.
This is the modulation spectrum, central in audio analysis.

We will see the same hierarchical structure in CNNs: first layer
detects edges; second layer detects combinations of edges; third layer
detects parts; eventually the network detects objects.  The CNN's
hierarchy is learned; the scattering transform's is mathematical.  They
are, structurally, the same hierarchy.

%------------------------------------------------------------------------
\section{Closing thoughts}
%------------------------------------------------------------------------

We have walked through three a-priori representations: the DFT
(maximum frequency resolution, no time resolution), the DWT (a
flexible time-frequency trade-off), and the Scattering Transform
(multi-layer, locally invariant, deformation-stable).  Each is a
specific choice of $\Phi(\xb) = f(\mathbf{B}\xb)$ in the framework of
the previous chapter, with $\mathbf{B}$ a fixed matrix designed by
mathematicians and $f$ either the identity, the magnitude
$\abs{\cdot}$, or a composition of those.

The three transforms answer different questions about a signal.  The
DFT answers ``what frequencies are present?''  The DWT answers ``what
frequencies are present, and when?''  The Scattering Transform
answers ``what frequencies are present, when, and how do they
interact?''  Each successive question is more demanding; each requires
more structure in the representation.  The progression suggests a
natural fourth step: ``what features are present, and we don't even
know what features we mean?''  When we don't know what kinds of
patterns to look for, we let the data tell us.  That is the subject of
the next chapter --- representations whose reference vectors are
\emph{learned}.

The progression also previews the architectural pattern we will
re-encounter many times: \emph{projection}, \emph{nonlinearity},
\emph{stacking}.  These three together make depth meaningful.  Without
projection there is nothing to compose; without nonlinearity composition
collapses; without stacking there is no hierarchy.  The Scattering
Transform demonstrates that all three can be in place without any
learning at all.  The point of learning, as we will see, is to let the
projections themselves adapt to whatever signals we throw at the
network --- to choose $\mathbf{B}$ from data rather than from a
mathematician's foresight.

% chapters/03_learned_representations.tex
%
% Part I, Chapter 3.  Learned representations: PCA, k-means, matching
% pursuit, neural networks (FC and conv).  Adapted from the geometric
% tour's Chapter 3 with matching pursuit added (it's central to Part III)
% and the "survey of architectures" section cut (deferred to Part II,
% which builds those architectures in depth, and the closing chapter of
% Part III, which explains what we deliberately didn't cover).

\chapter{Learned representations}
\labch{learned}

\epigraph{All models are wrong, but some are useful.}{George E.~P.~Box}

%------------------------------------------------------------------------
\section{When the basis comes from data}
%------------------------------------------------------------------------

\noindent The previous chapter built three representations whose
reference vectors --- sinusoids for Fourier, scaled and shifted wavelets
for the DWT, products of those for Scattering --- were chosen by
mathematicians in advance.  This works beautifully when we know what
property we want to expose.  We want pitch, so we project onto
sinusoids.  We want time-frequency structure, so we project onto
wavelets.  The chosen reference vectors carry our prior knowledge about
the signal.

\marginnote{%
\textbf{The bias-variance trade-off.}  Predefined bases are
\emph{low-variance}: they don't depend on which data you fed them, so
they don't overfit.  They're also \emph{high-bias}: if your signal
doesn't decompose well into the chosen basis, no amount of clever
analysis will help.  Learned bases reverse the trade-off.  This is the
classical bias-variance picture, made concrete.}

But what do we do when we don't know what to look for?  When we have a
million labeled images of cats and dogs and want a representation that
distinguishes them, we have no idea what mathematical function picks
out the relevant cat-vs-dog features.  When we have orchestral recordings
and want a representation that captures a composer's style, we cannot
write down by hand what reference vectors would expose ``Mahler-ness.''
The features are real but they are not in any catalogue we know how to
consult.

The natural move is to \emph{let the data choose the basis}.  Rather
than designing the reference vectors mathematically, we look at many
examples of signals and derive a set of vectors that, in some sense,
captures what is common or what is distinguishing.  This is what
\emph{learned representations} do.  The basis is no longer in a textbook;
it lives in a parameter array computed from data.

This chapter introduces four learned representations of increasing
sophistication.  Principal Component Analysis (Section~\ref{sec:pca})
finds the directions of greatest variance in a dataset.  $K$-means
(Section~\ref{sec:kmeans}) finds prototypes that group similar
examples.  Matching pursuit (Section~\ref{sec:mp}) builds a sparse,
interpretable approximation of a signal from a large dictionary of
candidates.  Neural networks (Section~\ref{sec:nn}) learn nonlinear,
multi-layered representations whose components are themselves
learned-from-data projections.

In all four cases, the underlying structure is the same as in the
previous chapter: project the signal onto reference vectors, possibly
apply a nonlinearity, possibly stack layers.  The only difference is
where the reference vectors come from.

%------------------------------------------------------------------------
\section{Principal Component Analysis}
\label{sec:pca}
%------------------------------------------------------------------------

PCA is the simplest learned representation: a linear projection whose
basis vectors are the directions of maximum variance in the data.

Given a data matrix $\mathbf{X} \in \R^{N \times d}$ with one example per
row (so $N$ examples of dimension $d$), PCA proceeds in three steps.

\paragraph{1. Centre the data.}  Subtract the per-feature mean so that
the columns of $\mathbf{X}$ have zero mean.  This is essential because
the geometric machinery of PCA assumes the data is centred at the
origin.

\paragraph{2. Compute the covariance.}  Form
\begin{equation}
    \mathbf{C} \;=\; \frac{1}{N}\, \mathbf{X}^\top \mathbf{X}.
\end{equation}
The $(i, j)$ entry of $\mathbf{C}$ is the covariance between feature $i$
and feature $j$.

\paragraph{3. Diagonalise.}  Find the eigenvalues $\lambda_1 \ge
\lambda_2 \ge \cdots \ge \lambda_d \ge 0$ and corresponding eigenvectors
$\mathbf{e}_1, \ldots, \mathbf{e}_d$ of $\mathbf{C}$.  These eigenvectors
are the \emph{principal components} --- the directions in feature space
along which the data has the most variance.  The first eigenvector
$\mathbf{e}_1$ is the single direction along which the data spreads
most; $\mathbf{e}_2$ is the direction of next greatest spread
perpendicular to $\mathbf{e}_1$; and so on.

To get a $k$-dimensional representation we project onto the first $k$
eigenvectors.  Let $\mathbf{E}_k \in \R^{d \times k}$ be the matrix
whose columns are $\mathbf{e}_1, \ldots, \mathbf{e}_k$.  Then
\begin{equation}
    \mathbf{Z} \;=\; \mathbf{X}\, \mathbf{E}_k
\end{equation}
is the data expressed in the $k$-dimensional principal-component basis.

\marginnote{%
PCA is mathematically a singular value decomposition (SVD) of the data
matrix $\mathbf{X}$ in disguise: the principal components are the right
singular vectors of $\mathbf{X}$.  When we work with PCA in Part~II,
the SVD is the practical computation.}

\paragraph{PCA in the framework.}  Compare $\mathbf{Z} = \mathbf{X}\,
\mathbf{E}_k$ with the analysis formula
$\mathbf{c} = \mathbf{B}\, \xb$ from the previous chapter.  PCA is
exactly that, with $\mathbf{B} = \mathbf{E}_k^\top$ (the principal-component
matrix) and the data being analysed simultaneously rather than one
example at a time.  The reference vectors are the eigenvectors.  The
only thing that distinguishes PCA from the Fourier transform, at the
level of mathematical structure, is that the basis vectors come from
the data instead of from $e^{i 2 \pi k n / N}$.

\paragraph{What PCA captures.}  PCA is a tool for finding low-dimensional
linear structure in high-dimensional data.  When applied to natural
images, its principal components often look like blurry blobs and
gradients --- the most-variance directions in the space of $32 \times 32$
RGB images.  When applied to speech features, the principal components
capture coarse vowel-vs-consonant structure.  When applied to text
embeddings, they capture broad semantic axes (topic, sentiment).

What PCA cannot capture is nonlinear structure.  The data may lie on
a curved manifold in feature space; PCA will fit it with a flat
hyperplane and lose all the curvature.  For nonlinear structure we
need either to engineer nonlinear features by hand (a hard problem) or
to learn them with a network (the rest of this chapter).

%------------------------------------------------------------------------
\section{$K$-means and prototypes}
\label{sec:kmeans}
%------------------------------------------------------------------------

$K$-means takes a different geometric view.  Rather than finding
directions of variance, it finds \emph{prototypes}: a small set of
representative points such that every other data point is close to one
of them.

\paragraph{The algorithm.}  Given data points $\xb_1, \ldots, \xb_N$
and a chosen number of clusters $K$:

\begin{enumerate}
    \item Initialise $K$ prototype vectors $\boldsymbol\mu_1, \ldots,
    \boldsymbol\mu_K$ (often by randomly picking $K$ of the data
    points).

    \item \emph{Assign} each data point $\xb_n$ to the nearest
    prototype:
    \[
        C_n \;=\; \arg\min_k \norm{\xb_n - \boldsymbol\mu_k}.
    \]

    \item \emph{Update} each prototype to the mean of the data points
    assigned to it:
    \[
        \boldsymbol\mu_k \;=\; \frac{1}{|\{n : C_n = k\}|}
        \sum_{n : C_n = k} \xb_n.
    \]

    \item Repeat steps 2--3 until the assignments stop changing.
\end{enumerate}

\paragraph{$K$-means in the framework.}  $K$-means looks like a
different kind of object from PCA --- a discrete clustering rather than
a linear projection.  But viewed through our framework it is still a
representation by reference vectors: every signal $\xb$ is summarised
by the index $C(\xb)$ of its nearest prototype, and the prototypes
$\{\boldsymbol\mu_k\}$ are the reference vectors.  The ``projection''
is the assignment operation, which is \emph{not} a linear map but a
piecewise-constant function of $\xb$.  This is our first nonlinear
representation, but the nonlinearity is hidden in the
\emph{arg-min}.

\marginnote{%
$K$-means and PCA are sometimes presented as alternatives, but they
answer different questions.  PCA: ``in what direction does the data
spread the most?''  $K$-means: ``what are the typical examples?''
A composer's archive analysed with PCA might give you a few axes of
stylistic variation.  Analysed with $K$-means, it might give you a
handful of representative pieces.  Both are useful; the right tool
depends on what you want to do.}

A useful matrix formulation: let $\mathbf{M} \in \R^{K \times d}$ be the
matrix whose rows are the prototypes, and let $\mathbf{C} \in \R^{N
\times K}$ be the binary assignment matrix with $C_{nk} = 1$ if example
$n$ is in cluster $k$ and zero otherwise.  Then the reconstruction
$\hat{\mathbf{X}} = \mathbf{C}\, \mathbf{M}$ replaces each data point
with its assigned prototype, and $K$-means minimises the reconstruction
error $\norm{\mathbf{X} - \mathbf{C}\, \mathbf{M}}_F^2$.  This is
structurally an extreme dimensionality reduction --- each example is
reduced to one integer in $\{1, \ldots, K\}$.

%------------------------------------------------------------------------
\section{Matching pursuit and sparse dictionaries}
\label{sec:mp}
%------------------------------------------------------------------------

The two representations so far --- PCA and $K$-means --- both insist on
a small basis: $k$ eigenvectors or $K$ prototypes.  Matching pursuit
takes the opposite tack: provide a \emph{large} dictionary of candidate
reference vectors, and let the analysis algorithm choose a small subset
of them per signal.  This is the right idea when no fixed-size basis
would do justice to the variety of signals we want to represent.

\paragraph{Setup.}  Let $\mathbf{D} \in \R^{d \times M}$ be a
\emph{dictionary}: a matrix whose $M$ columns $\mathbf{d}_1, \ldots,
\mathbf{d}_M$ are candidate reference vectors (atoms).  We allow $M$
to be much larger than $d$; the dictionary is \emph{overcomplete}.
Atoms could be Gabor wavelets at many scales and frequencies, samples
from a corpus of audio recordings, or anything else relevant.

The goal is to approximate a signal $\xb$ as a sparse linear combination
of atoms:
\begin{equation}
    \xb \;\approx\; \sum_{k=1}^{K} \alpha_k\, \mathbf{d}_{i_k},
    \qquad
    K \ll M.
\end{equation}
We want to find which $K$ atoms (the indices $i_1, \ldots, i_K$) and
which coefficients $\alpha_1, \ldots, \alpha_K$ minimise the
approximation error.  Choosing the best subset of $K$ atoms out of $M$
is combinatorially explosive --- there are $\binom{M}{K}$ subsets to
consider.  In practice we use a \emph{greedy} algorithm.

\paragraph{The matching pursuit algorithm.}  At each step we pick the
atom that explains the largest fraction of what remains.

\begin{enumerate}
    \item Start with residual $\mathbf{r}_0 = \xb$.
    \item For $k = 1, 2, \ldots, K$:
    \begin{enumerate}
        \item Find the atom most correlated with the current residual:
        \[
            i_k \;=\; \arg\max_i \; \abs{\inner{\mathbf{r}_{k-1}}{\mathbf{d}_i}}.
        \]
        \item Compute its coefficient: $\alpha_k = \inner{\mathbf{r}_{k-1}}{\mathbf{d}_{i_k}}$
        (assuming dictionary atoms are unit-norm).
        \item Subtract the projection: $\mathbf{r}_k = \mathbf{r}_{k-1}
        - \alpha_k \mathbf{d}_{i_k}$.
    \end{enumerate}
    \item Stop when $\norm{\mathbf{r}_K}$ is small enough, or when $K$
    reaches a budget.
\end{enumerate}

\noindent The algorithm produces an approximation $\hat\xb = \sum_k
\alpha_k \mathbf{d}_{i_k}$ that uses only $K$ of the $M$ available
atoms.

\paragraph{Why this is a representation.}  Matching pursuit produces,
for each input signal $\xb$, a sparse list of (atom, coefficient) pairs.
This list \emph{is} the representation.  Unlike a dense Fourier
spectrum or a flat PCA projection, matching pursuit returns only
those components that are large; the rest are implicit zeros.  The
representation is naturally \emph{interpretable}: each atom is a
specific pattern, and the coefficient tells you how much of that
pattern is in the signal.

\marginnote{%
Matching pursuit was introduced by Mallat and Zhang in 1993.  Its
descendants include orthogonal matching pursuit (OMP), basis pursuit
(which solves a related $L^1$ minimisation problem), and the broader
field of sparse coding and compressed sensing.  When the dictionary
itself is also learned from data, the algorithm is called
\emph{dictionary learning} or \emph{sparse coding}; this borders on
unsupervised neural-network methods.}

\paragraph{Matching pursuit in the framework.}  Matching pursuit is a
representation by reference vectors --- the atoms --- but the
``projection'' is no longer a single matrix multiplication.  It is the
output of a small algorithm (the greedy selection loop).  The basis is
also no longer a basis in the strict sense: the dictionary is redundant,
and we deliberately use only a few atoms per signal.

This breaks the clean ``matrix in, matrix out'' picture of the previous
chapter, but it gains two things in exchange.  First, \emph{sparsity}:
the representation has very few non-zero components, which is exactly
what we often want for compression and for interpretation.  Second,
\emph{adaptivity}: each signal gets its own custom subset of atoms,
matched to its content, rather than being squeezed through the same
projection as every other signal.

\paragraph{Why matching pursuit matters here.}  We will meet matching
pursuit again in Part~III, where it plays the role of a learned-but-not-neural
codec: a way to decompose orchestral sounds into a sequence of
atom-indices that captures their timbre and structure.  These
atom-indices then become the input to a transformer model that learns
to generate music.  In this sense matching pursuit gives us the best
of both worlds: an interpretable codec like Fourier, and a learned
basis like a neural network.

This pattern --- combine an a-priori transform with a learned model on
its output --- is increasingly common.  It is the basis of AudioLM,
MusicLM, and several other recent music-generation systems.  We will
build our own version in Part~III.

%------------------------------------------------------------------------
\section{Neural networks}
\label{sec:nn}
%------------------------------------------------------------------------

We come to the modern centre of mass.  A \emph{neural network} is a
representation built from many layers, each of which is a learned
linear projection followed by a nonlinearity.

\paragraph{A single layer.}  The basic building block, repeated many
times, is
\begin{equation}
    \zb \;=\; f(\Wb \xb + \mathbf{b}).
    \label{eq:nn_layer}
\end{equation}
Here $\Wb$ is a learned matrix of weights, $\mathbf{b}$ is a learned
bias vector that allows the projection to be affine rather than purely
linear, and $f$ is a fixed nonlinearity applied to each component (most
commonly ReLU, $f(x) = \max(0, x)$).

The inner product structure is right there in $\Wb \xb$.  Each row of
$\Wb$ is a reference vector $\wb_i$, and the $i$-th component of the
output is
\begin{equation}
    z_i \;=\; f\!\left( \inner{\wb_i}{\xb} + b_i \right).
\end{equation}
This is the inner-product analysis of the previous chapter, with two
modifications: the basis vectors $\wb_i$ are learned, and a nonlinearity
$f$ is applied to each result.

\paragraph{Logistic regression as a one-layer network.}  When the
nonlinearity is the sigmoid $\sigma(x) = 1 / (1 + e^{-x})$ and the
output is interpreted as a probability, equation~\eqref{eq:nn_layer} is
exactly \emph{logistic regression} --- the classical statistical model
for binary classification.  This is not a coincidence: a neural network
is, layer for layer, a stack of generalised regressions, each producing
features for the next.

\paragraph{Multiple layers.}  Stacking layers gives a deep network:
\begin{equation}
    \zb^{(L)} \;=\; f_L\!\left( \Wb_L \, f_{L-1}\!\left( \Wb_{L-1}\,
    \cdots f_1\!\left( \Wb_1 \xb + \mathbf{b}_1 \right) \cdots
    + \mathbf{b}_{L-1} \right) + \mathbf{b}_L \right).
    \label{eq:nn_deep}
\end{equation}
Without the nonlinearities $f_i$, this composition would collapse into
a single linear transformation --- the product $\Wb_L \cdots \Wb_1$ ---
giving us no more expressive power than a one-layer network.  The
nonlinearity is what makes depth meaningful.  We met this argument
already in the Scattering Transform, where the absolute value played
the same role.

\marginnote{%
\textbf{The universal approximation theorem.}  A sufficiently wide
two-layer network with a non-polynomial nonlinearity can approximate
any continuous function on a compact domain to any desired accuracy.
This is a celebrated theoretical result, and it sounds powerful, but
it understates the role of depth.  A two-layer network may need
exponentially many neurons to do what a deeper network does with
polynomially many.  Depth is not just possible; it is efficient.}

\paragraph{How the weights are learned.}  This is the algorithmic heart
of modern machine learning.  Given a dataset of input-output pairs
$\{(\xb_n, \yb_n)\}$, we define a \emph{loss function} $\mathcal{L}$
that measures how badly the network's predictions match the targets.
We then adjust the weights to reduce the loss by \emph{gradient
descent}:
\begin{equation}
    \Wb_\ell \;\leftarrow\; \Wb_\ell \,-\, \eta\,
    \frac{\partial \mathcal{L}}{\partial \Wb_\ell},
\end{equation}
with a learning rate $\eta$.  Computing the gradients
$\partial \mathcal{L} / \partial \Wb_\ell$ for every layer is done by
\emph{backpropagation}, which is the multivariate chain rule applied
methodically from the output layer back to the input layer.  Part~II
of this book derives backpropagation in detail, implements it in C++,
and verifies it with numerical gradient checks.

\paragraph{Encoders and decoders.}  We can use depth-with-nonlinearity
to build two complementary structures:

\begin{itemize}
    \item An \emph{encoder} starts with a high-dimensional input and
    produces a progressively lower-dimensional representation,
    distilling the input to a compact latent code.  Each layer has
    fewer outputs than inputs.
    \item A \emph{decoder} reverses the process: starts with a small
    latent vector and produces a high-dimensional output.  Each layer
    has more outputs than inputs.
\end{itemize}

\noindent An \emph{autoencoder} chains the two: encoder, latent
bottleneck, decoder, trained so the output reconstructs the input.
The latent bottleneck forces the network to find a compact
representation that loses as little information as possible.  When the
encoder and decoder are both deep nonlinear networks, the autoencoder
is a learned analogue of PCA --- with the same goal (dimensionality
reduction) but able to capture curved manifolds rather than just flat
hyperplanes.

A \emph{variational autoencoder} extends the autoencoder to a
generative model: the bottleneck is interpreted as a probability
distribution, and we can sample new outputs by sampling from the
latent distribution and running them through the decoder.  Part~II
builds both, and the audio VAE chapter of Part~II uses one to morph
between sounds.

%------------------------------------------------------------------------
\section{Convolutional neural networks}
%------------------------------------------------------------------------

A convolutional neural network applies the inner product locally, like
the Scattering Transform of the previous chapter, but with the filters
learned from data rather than designed by hand.  The architectural
template is exactly the same:

\begin{itemize}
    \item Convolve the input with a bank of filters.
    \item Apply a nonlinearity.
    \item Possibly pool (spatial downsampling).
    \item Repeat for several layers.
    \item Finish with a few dense layers for classification.
\end{itemize}

\noindent What changes from the scattering picture is that the filters
are no longer wavelets chosen for analytical convenience; they are
free parameters trained jointly with the rest of the network.  Each
filter is a small inner product applied at every location in the input:
the same inner product, slid across the signal.

\paragraph{Why share weights.}  The defining feature of a CNN compared
to a fully-connected network is \emph{weight sharing}: the same filter
weights are used at every spatial position.  This has three
consequences.  It reduces the number of parameters dramatically (a
$3 \times 3$ filter on an image of any size has only 9 weights, not 9
weights per image location).  It builds in \emph{translation
equivariance}: shifting the input shifts the output.  And it encodes
a prior of locality: each output value depends on a small local
neighbourhood, capturing the fact that pixel relationships are
usually nearby.

These three are not optional features; they are the architectural prior
that makes CNNs work on images and audio while fully-connected networks
struggle.

\marginnote{%
\textbf{Empirically interpretable.}  Train a CNN on natural images and
the first-layer filters --- which are easy to visualise as small RGB
patches --- come out looking like oriented edge detectors and
colour-opponent blobs.  These are exactly the kinds of features that
classical signal-processing theory tells us \emph{should} be useful for
natural images.  Krizhevsky's AlexNet paper (2012) and Zeiler and
Fergus's 2014 follow-up made this picture famous; we reproduce it in
Part~II's CIFAR-10 example.}

\paragraph{The hierarchy.}  In a multi-layer CNN, early layers learn
low-level features (edges, colour patches, texture primitives), middle
layers learn combinations (corners, simple shapes, repeating textures),
and late layers learn task-specific concepts (object parts, faces,
musical-instrument families).  This hierarchical organisation emerges
naturally from gradient descent on a classification loss; nothing in
the training procedure explicitly asks for it.

The same hierarchical pattern appeared in the Scattering Transform:
first-layer wavelet coefficients capture frequency content, second-layer
coefficients capture modulations of that content.  The CNN learns the
same kind of hierarchy, but discovers what the relevant features
\emph{are} from the data rather than receiving them as wavelets.  This
is the gain from learning, and it is the gain that explains why CNNs
dominate image and audio tasks in practice.

%------------------------------------------------------------------------
\section{The unified picture}
%------------------------------------------------------------------------

We have now seen, across two chapters, a sequence of representations:

\begin{itemize}
    \item \textbf{DFT}: linear projection onto sinusoids.
    \item \textbf{DWT}: linear projection onto scaled-shifted wavelets.
    \item \textbf{Scattering}: multi-layer wavelet cascade with absolute
    value as nonlinearity.
    \item \textbf{PCA}: linear projection onto eigenvectors of the data
    covariance.
    \item \textbf{$K$-means}: piecewise-constant assignment to learned
    prototypes.
    \item \textbf{Matching pursuit}: sparse selection from a learned (or
    designed) overcomplete dictionary.
    \item \textbf{Fully-connected network}: stacked nonlinear projections
    with learned dense weights.
    \item \textbf{Convolutional network}: stacked nonlinear projections
    with learned local weights, shared across spatial position.
\end{itemize}

\noindent Every one of these fits the template $\Phi(\xb) = f(\mathbf{B}
\xb)$, possibly with several layers stacked, possibly with the
projection $\mathbf{B}\xb$ replaced by a convolution.  The two
axes of variation are: (1) where the reference vectors come from
(predefined or learned), and (2) whether locality is enforced
(global inner products versus local convolutions).  This gives a $2
\times 2$ table:

\begin{center}
\begin{tabular}{l|cc}
        & \emph{Global} (inner product) & \emph{Local} (convolution) \\ \hline
\emph{Designed basis} & DFT, PCA, $K$-means & DWT, Scattering \\
\emph{Learned basis}  & FC network, matching pursuit & CNN \\
\end{tabular}
\end{center}

The progression of this book is along the diagonal: from designed-global
(Chapter~\ref{ch:apriori}) through designed-local and learned-global
(this chapter) to learned-local (the CNN, and the conv-VAE we will
build in Part~II).  Each step adds capability; each step also adds
opacity, because learned representations are harder to interpret than
designed ones.

The next chapter closes Part~I by working a concrete problem ---
musical style transfer --- through each cell of this table.  We will
see what \emph{the same problem looks like} when posed in the waveform
(Fourier), in the spectrum (DFT), in the wavelet domain (DWT), via
sparse atoms (matching pursuit), and via a learned latent space
(autoencoder).  The exercise will make vivid how the choice of
representation determines what is musically possible.

% \input{chapters/04_style_transfer_problem.tex}     % bridge to Part II


% =============================================================================
% PART II --- TECHNOLOGY
% =============================================================================
% Most chapters here are generated by pandoc from ../notes/*.md.
% See book/Makefile.  The \input lines are commented until the build pipeline
% has produced them.
% =============================================================================
\pagelayout{wide}
\addpart{Technology}
\pagelayout{margin}

% Part II chapters are generated from ../notes/*.md by `make notes`
% (see book/Makefile).  They are currently disabled because the notes
% contain Unicode math characters (·, →, ₁, etc.) inside code blocks
% that need a pre-processing pass before LaTeX renders them cleanly.
% See book/TODO_part2.md for the preprocessing plan.

\input{generated/00b_why_cpp.tex}
\input{generated/01_matrix.tex}
\input{generated/02_backprop.tex}
\input{generated/02b_one_weight.tex}
\input{generated/03_activations.tex}
\input{generated/03b_jacobian.tex}
\input{generated/04_cross_entropy.tex}
\input{generated/05_network.tex}
\input{generated/06_iris.tex}
\input{generated/07_optimizers.tex}
\input{generated/08_dropout.tex}
\input{generated/09_mnist.tex}
\input{generated/10_fft.tex}
\input{generated/11_audio_features.tex}
\input{generated/12_sound_classifier.tex}
\input{generated/13_autoencoder.tex}
\input{generated/14_vae.tex}
\input{generated/15_weight_visualization.tex}
\input{generated/16_audio_vae.tex}
\input{generated/17_performance.tex}
\input{generated/18_from_matrices_to_tensors.tex}
\input{generated/19_convnets.tex}


% =============================================================================
% PART III --- CREATION
% =============================================================================
\pagelayout{wide}
\addpart{Creation}
\pagelayout{margin}

% \input{chapters/20_transformations_without_learning.tex}
% \input{chapters/21_unsupervised_hybridizations.tex}
% \input{chapters/22_computer_assisted_orchestration.tex}
% \input{chapters/23_neural_design_for_music.tex}
% \input{chapters/24_a_philosophical_offering.tex}


% =============================================================================
% APPENDICES
% =============================================================================
\appendix
\pagelayout{wide}
\addpart{Appendices}
\pagelayout{margin}

% \input{chapters/A_matrix_theory.tex}
% \input{chapters/B_invariance_eigenvectors.tex}
% \input{chapters/C_complex_and_euler.tex}
% \input{chapters/D_convolution_theorem.tex}
% \input{chapters/E_chain_rule.tex}


%-------------------------------------------------------------------------------
\backmatter
\setchapterstyle{plain}

\defbibnote{bibnote}{References cited in this work.\par\bigskip}
\printbibliography[heading=bibintoc, title=Bibliography, prenote=bibnote]

\printindex

\end{document}
